% katibmath manual. Compile with XeLaTeX (or LuaLaTeX):  xelatex katibmath.tex
% The symbol table (katibmath-symbols.tex) is generated by
% scripts/build-ctan-package.ts in the الكاتب العلمي repository.
\documentclass[11pt]{article}
\usepackage[margin=2.2cm]{geometry}
\usepackage{katibmath}
\usepackage{longtable,booktabs}
\usepackage[hidelinks]{hyperref}
\setlength{\parindent}{0pt}
\setlength{\parskip}{0.5em}

\title{\textsf{katibmath}: Arabic mathematical notation\\for XeLaTeX and LuaLaTeX\\[0.3em]
  \large\textarabic{الترميز الرياضي العربي}}
\author{Muhammad Nabeel Shaesha}
\date{Version 0.1.0}

\begin{document}
\maketitle

\section{Introduction}

\textsf{katibmath} provides the Arabic mathematical notation used in Arabic
school and university texts: Arabic function names ($\jaa\arSeen$, $\nihaya_{\arSeen \to 0}$),
the Arabic Mathematical Alphabetic Symbols of Unicode ($\arDsTah$, $\arDsSad$, $\arDsHah$),
mirrored operators for Arabic-style formulas ($\rsum$, $\rint$, $\rsqrt{x}$),
the Arabic-script math signs ($\arCubeRoot$, $\arPercent$) and Arabic number sets.

It ships the OpenType math font \emph{Katib Math}: XITS Math with the Arabic
script of Amiri, under the SIL Open Font License. Every command in this package
is the one used by the Arabic writing app \emph{\textarabic{الكاتب العلمي}}
(Scientific Writer), so documents exported from the app compile unchanged.

\section{Usage}

\begin{verbatim}
\usepackage{katibmath}            % Katib Math becomes the math font
\usepackage[nomathfont]{katibmath} % keep your math font
\end{verbatim}

The package needs a Unicode engine (XeLaTeX or LuaLaTeX) and loads
\textsf{amsmath}, \textsf{graphicx}, \textsf{fontspec} and \textsf{unicode-math}.

With \textsf{polyglossia} or \textsf{babel} loaded for Arabic, their
\verb|\textarabic| typesets the Arabic words in formulas. Without them,
\textsf{katibmath} defines \verb|\textarabic| itself, using the Arabic of Katib
Math, so the commands also work in English documents.

\subsection*{Example}

\begin{verbatim}
\[ \nihaya_{\arSeen \to 0} \frac{\jaa\arSeen}{\arSeen} = 1 \qquad
   \rsum_{\arKaf=1}^{\arNoon} \arKaf \qquad
   \rint_0^1 f \qquad \rsqrt{\arSeen} \qquad \arReals \]
\end{verbatim}
\[ \nihaya_{\arSeen \to 0} \frac{\jaa\arSeen}{\arSeen} = 1 \qquad
   \rsum_{\arKaf=1}^{\arNoon} \arKaf \qquad
   \rint_0^1 f \qquad \rsqrt{\arSeen} \qquad \arReals \]

Bare Arabic letters typed in math are not shaped by the math font: use the
alphabet commands (\verb|\arSeen|, \verb|\arKaf| … in section~\ref{sec:symbols})
or wrap words in \verb|\textarabic{...}|, which works in math mode too.

\subsection*{Helpers}

\begin{description}
  \item[\texttt{\textbackslash katibmirror\{x\}}] $x$ flipped horizontally in the
    current math style: \verb|\katibmirror{\sum}| gives $\katibmirror{\sum}$.
  \item[\texttt{\textbackslash katibglyphs}] Katib Math as a text font.
  \item[\texttt{\textbackslash katibarabic}] Katib Math's Arabic, shaped.
\end{description}

\subsection*{Limits}

Right-to-left formulas (the whole formula laid out from right to left) are not
supported by \TeX's math engine; \textsf{katibmath} gives Arabic \emph{symbols}
inside left-to-right formulas.

\section{Symbols}\label{sec:symbols}
% Generated by scripts/build-ctan-package.ts. Do not edit by hand.
\subsection*{Functions and limits \quad \textarabic{الدوال والنهايات}}
\begin{longtable}{@{}l c c r@{}}
\toprule Command & Symbol & Example & \textarabic{الاسم} \\ \midrule \endhead
\verb|\jaa| & $\jaa$ & $\displaystyle \jaa$ & \textarabic{جيب الزاوية} \\
\verb|\jataa| & $\jataa$ & $\displaystyle \jataa$ & \textarabic{جيب تمام الزاوية} \\
\verb|\dhaa| & $\dhaa$ & $\displaystyle \dhaa$ & \textarabic{ظل الزاوية} \\
\verb|\dhataa| & $\dhataa$ & $\displaystyle \dhataa$ & \textarabic{ظل تمام الزاوية} \\
\verb|\qaa| & $\qaa$ & $\displaystyle \qaa$ & \textarabic{قاطع الزاوية} \\
\verb|\qataa| & $\qataa$ & $\displaystyle \qataa$ & \textarabic{قاطع تمام الزاوية} \\
\verb|\jaainv| & $\jaainv$ & $\displaystyle \jaainv$ & \textarabic{معكوس الجيب} \\
\verb|\jataainv| & $\jataainv$ & $\displaystyle \jataainv$ & \textarabic{معكوس جيب التمام} \\
\verb|\dhaainv| & $\dhaainv$ & $\displaystyle \dhaainv$ & \textarabic{معكوس الظل} \\
\verb|\jaaz| & $\jaaz$ & $\displaystyle \jaaz$ & \textarabic{الجيب الزائدي} \\
\verb|\jataaz| & $\jataaz$ & $\displaystyle \jataaz$ & \textarabic{جيب التمام الزائدي} \\
\verb|\dhaaz| & $\dhaaz$ & $\displaystyle \dhaaz$ & \textarabic{الظل الزائدي} \\
\verb|\lugh| & $\lugh$ & $\displaystyle \lugh$ & \textarabic{اللوغاريتم} \\
\verb|\lugt| & $\lugt$ & $\displaystyle \lugt$ & \textarabic{اللوغاريتم الطبيعي} \\
\verb|\nihaya| & $\nihaya$ & $\displaystyle \nihaya$ & \textarabic{النهاية} \\
\verb|\majmu| & $\majmu$ & $\displaystyle \majmu$ & \textarabic{المجموع} \\
\bottomrule
\end{longtable}
\subsection*{Arabic mathematical alphabet \quad \textarabic{الأبجدية الرياضية العربية}}
\begin{longtable}{@{}l c c r@{}}
\toprule Command & Symbol & Example & \textarabic{الاسم} \\ \midrule \endhead
\verb|\arAlef| & $\arAlef$ & $\displaystyle \arAlef$ & \textarabic{ا} \\
\verb|\arBeh| & $\arBeh$ & $\displaystyle \arBeh$ & \textarabic{ب} \\
\verb|\arJeem| & $\arJeem$ & $\displaystyle \arJeem$ & \textarabic{ج} \\
\verb|\arDal| & $\arDal$ & $\displaystyle \arDal$ & \textarabic{د} \\
\verb|\arWaw| & $\arWaw$ & $\displaystyle \arWaw$ & \textarabic{و} \\
\verb|\arZain| & $\arZain$ & $\displaystyle \arZain$ & \textarabic{ز} \\
\verb|\arHah| & $\arHah$ & $\displaystyle \arHah$ & \textarabic{ح} \\
\verb|\arTah| & $\arTah$ & $\displaystyle \arTah$ & \textarabic{ط} \\
\verb|\arYeh| & $\arYeh$ & $\displaystyle \arYeh$ & \textarabic{ي} \\
\verb|\arKaf| & $\arKaf$ & $\displaystyle \arKaf$ & \textarabic{ك} \\
\verb|\arLam| & $\arLam$ & $\displaystyle \arLam$ & \textarabic{ل} \\
\verb|\arMeem| & $\arMeem$ & $\displaystyle \arMeem$ & \textarabic{م} \\
\verb|\arNoon| & $\arNoon$ & $\displaystyle \arNoon$ & \textarabic{ن} \\
\verb|\arSeen| & $\arSeen$ & $\displaystyle \arSeen$ & \textarabic{س} \\
\verb|\arAin| & $\arAin$ & $\displaystyle \arAin$ & \textarabic{ع} \\
\verb|\arFeh| & $\arFeh$ & $\displaystyle \arFeh$ & \textarabic{ف} \\
\verb|\arSad| & $\arSad$ & $\displaystyle \arSad$ & \textarabic{ص} \\
\verb|\arQaf| & $\arQaf$ & $\displaystyle \arQaf$ & \textarabic{ق} \\
\verb|\arReh| & $\arReh$ & $\displaystyle \arReh$ & \textarabic{ر} \\
\verb|\arSheen| & $\arSheen$ & $\displaystyle \arSheen$ & \textarabic{ش} \\
\verb|\arTeh| & $\arTeh$ & $\displaystyle \arTeh$ & \textarabic{ت} \\
\verb|\arTheh| & $\arTheh$ & $\displaystyle \arTheh$ & \textarabic{ث} \\
\verb|\arKhah| & $\arKhah$ & $\displaystyle \arKhah$ & \textarabic{خ} \\
\verb|\arThal| & $\arThal$ & $\displaystyle \arThal$ & \textarabic{ذ} \\
\verb|\arDad| & $\arDad$ & $\displaystyle \arDad$ & \textarabic{ض} \\
\verb|\arZah| & $\arZah$ & $\displaystyle \arZah$ & \textarabic{ظ} \\
\verb|\arGhain| & $\arGhain$ & $\displaystyle \arGhain$ & \textarabic{غ} \\
\verb|\arDotlessBeh| & $\arDotlessBeh$ & $\displaystyle \arDotlessBeh$ & \textarabic{ٮ بلا نقط} \\
\verb|\arDotlessNoon| & $\arDotlessNoon$ & $\displaystyle \arDotlessNoon$ & \textarabic{ن بلا نقط} \\
\verb|\arDotlessFeh| & $\arDotlessFeh$ & $\displaystyle \arDotlessFeh$ & \textarabic{ڡ بلا نقط} \\
\verb|\arDotlessQaf| & $\arDotlessQaf$ & $\displaystyle \arDotlessQaf$ & \textarabic{ٯ بلا نقط} \\
\verb|\arInitBeh| & $\arInitBeh$ & $\displaystyle \arInitBeh$ & \textarabic{ب مستهلّة} \\
\verb|\arInitJeem| & $\arInitJeem$ & $\displaystyle \arInitJeem$ & \textarabic{ج مستهلّة} \\
\verb|\arInitHeh| & $\arInitHeh$ & $\displaystyle \arInitHeh$ & \textarabic{ه مستهلّة} \\
\verb|\arInitHah| & $\arInitHah$ & $\displaystyle \arInitHah$ & \textarabic{ح مستهلّة} \\
\verb|\arInitYeh| & $\arInitYeh$ & $\displaystyle \arInitYeh$ & \textarabic{ي مستهلّة} \\
\verb|\arInitKaf| & $\arInitKaf$ & $\displaystyle \arInitKaf$ & \textarabic{ك مستهلّة} \\
\verb|\arInitLam| & $\arInitLam$ & $\displaystyle \arInitLam$ & \textarabic{ل مستهلّة} \\
\verb|\arInitMeem| & $\arInitMeem$ & $\displaystyle \arInitMeem$ & \textarabic{م مستهلّة} \\
\verb|\arInitNoon| & $\arInitNoon$ & $\displaystyle \arInitNoon$ & \textarabic{ن مستهلّة} \\
\verb|\arInitSeen| & $\arInitSeen$ & $\displaystyle \arInitSeen$ & \textarabic{س مستهلّة} \\
\verb|\arInitAin| & $\arInitAin$ & $\displaystyle \arInitAin$ & \textarabic{ع مستهلّة} \\
\verb|\arInitFeh| & $\arInitFeh$ & $\displaystyle \arInitFeh$ & \textarabic{ف مستهلّة} \\
\verb|\arInitSad| & $\arInitSad$ & $\displaystyle \arInitSad$ & \textarabic{ص مستهلّة} \\
\verb|\arInitQaf| & $\arInitQaf$ & $\displaystyle \arInitQaf$ & \textarabic{ق مستهلّة} \\
\verb|\arInitSheen| & $\arInitSheen$ & $\displaystyle \arInitSheen$ & \textarabic{ش مستهلّة} \\
\verb|\arInitTeh| & $\arInitTeh$ & $\displaystyle \arInitTeh$ & \textarabic{ت مستهلّة} \\
\verb|\arInitTheh| & $\arInitTheh$ & $\displaystyle \arInitTheh$ & \textarabic{ث مستهلّة} \\
\verb|\arInitKhah| & $\arInitKhah$ & $\displaystyle \arInitKhah$ & \textarabic{خ مستهلّة} \\
\verb|\arInitDad| & $\arInitDad$ & $\displaystyle \arInitDad$ & \textarabic{ض مستهلّة} \\
\verb|\arInitGhain| & $\arInitGhain$ & $\displaystyle \arInitGhain$ & \textarabic{غ مستهلّة} \\
\verb|\arTailJeem| & $\arTailJeem$ & $\displaystyle \arTailJeem$ & \textarabic{ج بذيل} \\
\verb|\arTailHah| & $\arTailHah$ & $\displaystyle \arTailHah$ & \textarabic{ح بذيل} \\
\verb|\arTailYeh| & $\arTailYeh$ & $\displaystyle \arTailYeh$ & \textarabic{ي بذيل} \\
\verb|\arTailLam| & $\arTailLam$ & $\displaystyle \arTailLam$ & \textarabic{ل بذيل} \\
\verb|\arTailNoon| & $\arTailNoon$ & $\displaystyle \arTailNoon$ & \textarabic{ن بذيل} \\
\verb|\arTailSeen| & $\arTailSeen$ & $\displaystyle \arTailSeen$ & \textarabic{س بذيل} \\
\verb|\arTailAin| & $\arTailAin$ & $\displaystyle \arTailAin$ & \textarabic{ع بذيل} \\
\verb|\arTailSad| & $\arTailSad$ & $\displaystyle \arTailSad$ & \textarabic{ص بذيل} \\
\verb|\arTailQaf| & $\arTailQaf$ & $\displaystyle \arTailQaf$ & \textarabic{ق بذيل} \\
\verb|\arTailSheen| & $\arTailSheen$ & $\displaystyle \arTailSheen$ & \textarabic{ش بذيل} \\
\verb|\arTailKhah| & $\arTailKhah$ & $\displaystyle \arTailKhah$ & \textarabic{خ بذيل} \\
\verb|\arTailDad| & $\arTailDad$ & $\displaystyle \arTailDad$ & \textarabic{ض بذيل} \\
\verb|\arTailGhain| & $\arTailGhain$ & $\displaystyle \arTailGhain$ & \textarabic{غ بذيل} \\
\verb|\arTailDotlessNoon| & $\arTailDotlessNoon$ & $\displaystyle \arTailDotlessNoon$ & \textarabic{DotlessNoon بذيل} \\
\verb|\arTailDotlessQaf| & $\arTailDotlessQaf$ & $\displaystyle \arTailDotlessQaf$ & \textarabic{ٯ بذيل} \\
\verb|\arStretchBeh| & $\arStretchBeh$ & $\displaystyle \arStretchBeh$ & \textarabic{ب ممدودة} \\
\verb|\arStretchJeem| & $\arStretchJeem$ & $\displaystyle \arStretchJeem$ & \textarabic{ج ممدودة} \\
\verb|\arStretchHeh| & $\arStretchHeh$ & $\displaystyle \arStretchHeh$ & \textarabic{ه ممدودة} \\
\verb|\arStretchHah| & $\arStretchHah$ & $\displaystyle \arStretchHah$ & \textarabic{ح ممدودة} \\
\verb|\arStretchTah| & $\arStretchTah$ & $\displaystyle \arStretchTah$ & \textarabic{ط ممدودة} \\
\verb|\arStretchYeh| & $\arStretchYeh$ & $\displaystyle \arStretchYeh$ & \textarabic{ي ممدودة} \\
\verb|\arStretchKaf| & $\arStretchKaf$ & $\displaystyle \arStretchKaf$ & \textarabic{ك ممدودة} \\
\verb|\arStretchMeem| & $\arStretchMeem$ & $\displaystyle \arStretchMeem$ & \textarabic{م ممدودة} \\
\verb|\arStretchNoon| & $\arStretchNoon$ & $\displaystyle \arStretchNoon$ & \textarabic{ن ممدودة} \\
\verb|\arStretchSeen| & $\arStretchSeen$ & $\displaystyle \arStretchSeen$ & \textarabic{س ممدودة} \\
\verb|\arStretchAin| & $\arStretchAin$ & $\displaystyle \arStretchAin$ & \textarabic{ع ممدودة} \\
\verb|\arStretchFeh| & $\arStretchFeh$ & $\displaystyle \arStretchFeh$ & \textarabic{ف ممدودة} \\
\verb|\arStretchSad| & $\arStretchSad$ & $\displaystyle \arStretchSad$ & \textarabic{ص ممدودة} \\
\verb|\arStretchQaf| & $\arStretchQaf$ & $\displaystyle \arStretchQaf$ & \textarabic{ق ممدودة} \\
\verb|\arStretchSheen| & $\arStretchSheen$ & $\displaystyle \arStretchSheen$ & \textarabic{ش ممدودة} \\
\verb|\arStretchTeh| & $\arStretchTeh$ & $\displaystyle \arStretchTeh$ & \textarabic{ت ممدودة} \\
\verb|\arStretchTheh| & $\arStretchTheh$ & $\displaystyle \arStretchTheh$ & \textarabic{ث ممدودة} \\
\verb|\arStretchKhah| & $\arStretchKhah$ & $\displaystyle \arStretchKhah$ & \textarabic{خ ممدودة} \\
\verb|\arStretchDad| & $\arStretchDad$ & $\displaystyle \arStretchDad$ & \textarabic{ض ممدودة} \\
\verb|\arStretchZah| & $\arStretchZah$ & $\displaystyle \arStretchZah$ & \textarabic{ظ ممدودة} \\
\verb|\arStretchGhain| & $\arStretchGhain$ & $\displaystyle \arStretchGhain$ & \textarabic{غ ممدودة} \\
\verb|\arStretchDotlessBeh| & $\arStretchDotlessBeh$ & $\displaystyle \arStretchDotlessBeh$ & \textarabic{ٮ ممدودة} \\
\verb|\arStretchDotlessFeh| & $\arStretchDotlessFeh$ & $\displaystyle \arStretchDotlessFeh$ & \textarabic{ڡ ممدودة} \\
\verb|\arLoopAlef| & $\arLoopAlef$ & $\displaystyle \arLoopAlef$ & \textarabic{ا معقودة} \\
\verb|\arLoopBeh| & $\arLoopBeh$ & $\displaystyle \arLoopBeh$ & \textarabic{ب معقودة} \\
\verb|\arLoopJeem| & $\arLoopJeem$ & $\displaystyle \arLoopJeem$ & \textarabic{ج معقودة} \\
\verb|\arLoopDal| & $\arLoopDal$ & $\displaystyle \arLoopDal$ & \textarabic{د معقودة} \\
\verb|\arLoopHeh| & $\arLoopHeh$ & $\displaystyle \arLoopHeh$ & \textarabic{ه معقودة} \\
\verb|\arLoopWaw| & $\arLoopWaw$ & $\displaystyle \arLoopWaw$ & \textarabic{و معقودة} \\
\verb|\arLoopZain| & $\arLoopZain$ & $\displaystyle \arLoopZain$ & \textarabic{ز معقودة} \\
\verb|\arLoopHah| & $\arLoopHah$ & $\displaystyle \arLoopHah$ & \textarabic{ح معقودة} \\
\verb|\arLoopTah| & $\arLoopTah$ & $\displaystyle \arLoopTah$ & \textarabic{ط معقودة} \\
\verb|\arLoopYeh| & $\arLoopYeh$ & $\displaystyle \arLoopYeh$ & \textarabic{ي معقودة} \\
\verb|\arLoopLam| & $\arLoopLam$ & $\displaystyle \arLoopLam$ & \textarabic{ل معقودة} \\
\verb|\arLoopMeem| & $\arLoopMeem$ & $\displaystyle \arLoopMeem$ & \textarabic{م معقودة} \\
\verb|\arLoopNoon| & $\arLoopNoon$ & $\displaystyle \arLoopNoon$ & \textarabic{ن معقودة} \\
\verb|\arLoopSeen| & $\arLoopSeen$ & $\displaystyle \arLoopSeen$ & \textarabic{س معقودة} \\
\verb|\arLoopAin| & $\arLoopAin$ & $\displaystyle \arLoopAin$ & \textarabic{ع معقودة} \\
\verb|\arLoopFeh| & $\arLoopFeh$ & $\displaystyle \arLoopFeh$ & \textarabic{ف معقودة} \\
\verb|\arLoopSad| & $\arLoopSad$ & $\displaystyle \arLoopSad$ & \textarabic{ص معقودة} \\
\verb|\arLoopQaf| & $\arLoopQaf$ & $\displaystyle \arLoopQaf$ & \textarabic{ق معقودة} \\
\verb|\arLoopReh| & $\arLoopReh$ & $\displaystyle \arLoopReh$ & \textarabic{ر معقودة} \\
\verb|\arLoopSheen| & $\arLoopSheen$ & $\displaystyle \arLoopSheen$ & \textarabic{ش معقودة} \\
\verb|\arLoopTeh| & $\arLoopTeh$ & $\displaystyle \arLoopTeh$ & \textarabic{ت معقودة} \\
\verb|\arLoopTheh| & $\arLoopTheh$ & $\displaystyle \arLoopTheh$ & \textarabic{ث معقودة} \\
\verb|\arLoopKhah| & $\arLoopKhah$ & $\displaystyle \arLoopKhah$ & \textarabic{خ معقودة} \\
\verb|\arLoopThal| & $\arLoopThal$ & $\displaystyle \arLoopThal$ & \textarabic{ذ معقودة} \\
\verb|\arLoopDad| & $\arLoopDad$ & $\displaystyle \arLoopDad$ & \textarabic{ض معقودة} \\
\verb|\arLoopZah| & $\arLoopZah$ & $\displaystyle \arLoopZah$ & \textarabic{ظ معقودة} \\
\verb|\arLoopGhain| & $\arLoopGhain$ & $\displaystyle \arLoopGhain$ & \textarabic{غ معقودة} \\
\verb|\arDsBeh| & $\arDsBeh$ & $\displaystyle \arDsBeh$ & \textarabic{ب مزدوجة} \\
\verb|\arDsJeem| & $\arDsJeem$ & $\displaystyle \arDsJeem$ & \textarabic{ج مزدوجة} \\
\verb|\arDsDal| & $\arDsDal$ & $\displaystyle \arDsDal$ & \textarabic{د مزدوجة} \\
\verb|\arDsWaw| & $\arDsWaw$ & $\displaystyle \arDsWaw$ & \textarabic{و مزدوجة} \\
\verb|\arDsZain| & $\arDsZain$ & $\displaystyle \arDsZain$ & \textarabic{ز مزدوجة} \\
\verb|\arDsHah| & $\arDsHah$ & $\displaystyle \arDsHah$ & \textarabic{ح مزدوجة} \\
\verb|\arDsTah| & $\arDsTah$ & $\displaystyle \arDsTah$ & \textarabic{ط مزدوجة} \\
\verb|\arDsYeh| & $\arDsYeh$ & $\displaystyle \arDsYeh$ & \textarabic{ي مزدوجة} \\
\verb|\arDsLam| & $\arDsLam$ & $\displaystyle \arDsLam$ & \textarabic{ل مزدوجة} \\
\verb|\arDsMeem| & $\arDsMeem$ & $\displaystyle \arDsMeem$ & \textarabic{م مزدوجة} \\
\verb|\arDsNoon| & $\arDsNoon$ & $\displaystyle \arDsNoon$ & \textarabic{ن مزدوجة} \\
\verb|\arDsSeen| & $\arDsSeen$ & $\displaystyle \arDsSeen$ & \textarabic{س مزدوجة} \\
\verb|\arDsAin| & $\arDsAin$ & $\displaystyle \arDsAin$ & \textarabic{ع مزدوجة} \\
\verb|\arDsFeh| & $\arDsFeh$ & $\displaystyle \arDsFeh$ & \textarabic{ف مزدوجة} \\
\verb|\arDsSad| & $\arDsSad$ & $\displaystyle \arDsSad$ & \textarabic{ص مزدوجة} \\
\verb|\arDsQaf| & $\arDsQaf$ & $\displaystyle \arDsQaf$ & \textarabic{ق مزدوجة} \\
\verb|\arDsReh| & $\arDsReh$ & $\displaystyle \arDsReh$ & \textarabic{ر مزدوجة} \\
\verb|\arDsSheen| & $\arDsSheen$ & $\displaystyle \arDsSheen$ & \textarabic{ش مزدوجة} \\
\verb|\arDsTeh| & $\arDsTeh$ & $\displaystyle \arDsTeh$ & \textarabic{ت مزدوجة} \\
\verb|\arDsTheh| & $\arDsTheh$ & $\displaystyle \arDsTheh$ & \textarabic{ث مزدوجة} \\
\verb|\arDsKhah| & $\arDsKhah$ & $\displaystyle \arDsKhah$ & \textarabic{خ مزدوجة} \\
\verb|\arDsThal| & $\arDsThal$ & $\displaystyle \arDsThal$ & \textarabic{ذ مزدوجة} \\
\verb|\arDsDad| & $\arDsDad$ & $\displaystyle \arDsDad$ & \textarabic{ض مزدوجة} \\
\verb|\arDsZah| & $\arDsZah$ & $\displaystyle \arDsZah$ & \textarabic{ظ مزدوجة} \\
\verb|\arDsGhain| & $\arDsGhain$ & $\displaystyle \arDsGhain$ & \textarabic{غ مزدوجة} \\
\verb|\arOpMeemHah| & $\arOpMeemHah$ & $\displaystyle \arOpMeemHah$ & \textarabic{مؤثر مح (مجموع)} \\
\verb|\arOpHahDal| & $\arOpHahDal$ & $\displaystyle \arOpHahDal$ & \textarabic{مؤثر حد (نهاية)} \\
\bottomrule
\end{longtable}
\subsection*{Mirrored operators \quad \textarabic{المؤثرات المعكوسة}}
\begin{longtable}{@{}l c c r@{}}
\toprule Command & Symbol & Example & \textarabic{الاسم} \\ \midrule \endhead
\verb|\rsum| & $\rsum$ & $\displaystyle \rsum$ & \textarabic{مجموع معكوس} \\
\verb|\rint| & $\rint$ & $\displaystyle \rint$ & \textarabic{تكامل معكوس} \\
\verb|\riint| & $\riint$ & $\displaystyle \riint$ & \textarabic{تكامل ثنائي معكوس} \\
\verb|\riiint| & $\riiint$ & $\displaystyle \riiint$ & \textarabic{تكامل ثلاثي معكوس} \\
\verb|\roint| & $\roint$ & $\displaystyle \roint$ & \textarabic{تكامل خطي مغلق معكوس} \\
\verb|\roiint| & $\roiint$ & $\displaystyle \roiint$ & \textarabic{تكامل سطحي مغلق معكوس} \\
\verb|\rsqrt{x}| & $\rsqrt{x}$ & $\displaystyle \rsqrt{x}$ & \textarabic{جذر معكوس} \\
\verb|\rroot{3}{x}| & $\rroot{3}{x}$ & $\displaystyle \rroot{3}{x}$ & \textarabic{جذر نوني معكوس} \\
\verb|\rpartial| & $\rpartial$ & $\displaystyle \rpartial$ & \textarabic{مشتقة جزئية معكوسة} \\
\verb|\rexists| & $\rexists$ & $\displaystyle \rexists$ & \textarabic{يوجد (معكوس)} \\
\verb|\rnexists| & $\rnexists$ & $\displaystyle \rnexists$ & \textarabic{لا يوجد (معكوس)} \\
\verb|\ranglesign| & $\ranglesign$ & $\displaystyle \ranglesign$ & \textarabic{زاوية معكوسة} \\
\verb|\rmeasuredangle| & $\rmeasuredangle$ & $\displaystyle \rmeasuredangle$ & \textarabic{زاوية مقيسة معكوسة} \\
\verb|\rpropto| & $\rpropto$ & $\displaystyle \rpropto$ & \textarabic{يتناسب مع (معكوس)} \\
\verb|\rcomplement| & $\rcomplement$ & $\displaystyle \rcomplement$ & \textarabic{متمّمة معكوسة} \\
\verb|\arGcd| & $\arGcd$ & $\displaystyle \arGcd$ & \textarabic{القاسم المشترك الأكبر (ق.م.أ)} \\
\verb|\arLcm| & $\arLcm$ & $\displaystyle \arLcm$ & \textarabic{المضاعف المشترك الأصغر (م.م.أ)} \\
\verb|\arDet| & $\arDet$ & $\displaystyle \arDet$ & \textarabic{المحدد} \\
\verb|\arFact{x}| & $\arFact{x}$ & $\displaystyle \arFact{x}$ & \textarabic{المضروب (علامة التعجب قبل العدد)} \\
\bottomrule
\end{longtable}
\subsection*{Signs and number sets \quad \textarabic{إشارات ومجموعات أعداد}}
\begin{longtable}{@{}l c c r@{}}
\toprule Command & Symbol & Example & \textarabic{الاسم} \\ \midrule \endhead
\verb|\arCubeRoot| & $\arCubeRoot$ & $\displaystyle \arCubeRoot$ & \textarabic{جذر تكعيبي عربي ؆} \\
\verb|\arFourthRoot| & $\arFourthRoot$ & $\displaystyle \arFourthRoot$ & \textarabic{جذر رابع عربي ؇} \\
\verb|\arRay| & $\arRay$ & $\displaystyle \arRay$ & \textarabic{شعاع ؈} \\
\verb|\arPercent| & $\arPercent$ & $\displaystyle \arPercent$ & \textarabic{نسبة مئوية ٪} \\
\verb|\arPermille| & $\arPermille$ & $\displaystyle \arPermille$ & \textarabic{في الألف ؉} \\
\verb|\arPerTenThousand| & $\arPerTenThousand$ & $\displaystyle \arPerTenThousand$ & \textarabic{في العشرة آلاف ؊} \\
\verb|\arNaturals| & $\arNaturals$ & $\displaystyle \arNaturals$ & \textarabic{الأعداد الطبيعية ط} \\
\verb|\arIntegers| & $\arIntegers$ & $\displaystyle \arIntegers$ & \textarabic{الأعداد الصحيحة ص} \\
\verb|\arRationals| & $\arRationals$ & $\displaystyle \arRationals$ & \textarabic{الأعداد النسبية ن} \\
\verb|\arReals| & $\arReals$ & $\displaystyle \arReals$ & \textarabic{الأعداد الحقيقية ح} \\
\verb|\arComplex| & $\arComplex$ & $\displaystyle \arComplex$ & \textarabic{الأعداد المركّبة ك} \\
\bottomrule
\end{longtable}
\subsection*{Relations and notation \quad \textarabic{العلاقات والترميز}}
\begin{longtable}{@{}l c c r@{}}
\toprule Command & Symbol & Example & \textarabic{الاسم} \\ \midrule \endhead
\verb|\arMod| & $\arMod$ & $\displaystyle \arMod$ & \textarabic{باقي القسمة (مودولو)} \\
\bottomrule
\end{longtable}


\section{Licence}

The package code is distributed under the \LaTeX\ Project Public License 1.3c or
later. The font \texttt{KatibMath-Regular.otf} is distributed under the SIL Open
Font License 1.1 (\texttt{OFL.txt}); it is derived from XITS Math and Amiri.

\end{document}
